\documentclass[10pt,letterpaper]{article}

\usepackage[margin=0.42in]{geometry}
\usepackage[T1]{fontenc}
\usepackage[utf8]{inputenc}
\usepackage{enumitem}
\usepackage[hidelinks]{hyperref}
\usepackage{titlesec}
\usepackage{xcolor}

\pagestyle{empty}
\setlength{\parindent}{0pt}
\setlength{\parskip}{0pt}
\setlist[itemize]{leftmargin=*, itemsep=0pt, topsep=1pt}
\titleformat{\section}{\large\bfseries}{}{0pt}{}[\titlerule]
\titlespacing*{\section}{0pt}{5pt}{3pt}

\newcommand{\resumeItem}[1]{\item #1}
\newcommand{\resumeSubheading}[4]{
  \textbf{#1} \hfill #2 \\
  \textit{#3} \hfill \textit{#4}
}
\newcommand{\projectHeading}[2]{
  \textbf{#1} \hfill \textit{#2}
}

\begin{document}

\begin{center}
  {\LARGE \textbf{Beiyan Liu}} \\
  \vspace{2pt}
  Agent Systems Engineer \,|\, Context Infrastructure \,|\, Agent Harness \,|\, Evaluation \\
  \vspace{2pt}
  \href{mailto:YOUR_EMAIL@example.com}{YOUR\_EMAIL@example.com} \,|\,
  \href{https://github.com/Barytes}{github.com/Barytes} \,|\,
  China / Remote
\end{center}

\section{Profile}
Engineer and researcher building agent systems around context, knowledge workbenches, tool use, and evaluation-first reliability. Public projects focus on local-first knowledge bases, shared research context layers, and minimal agent harnesses where context is explicit, tool use is traceable, and failures can become regression cases.

\section{Selected Projects}

\projectHeading{\href{https://github.com/Barytes/gogo}{gogo} -- Local LLM-Wiki Desktop Workbench}{Open source / released prototype}
\begin{itemize}
  \resumeItem{Built a local-first desktop workbench for \texttt{llm-wiki}-style knowledge bases, combining \texttt{Wiki / Chat} modes, local \texttt{wiki / raw / inbox} browsing, Pi agent integration, session history, knowledge-base switching, and onboarding.}
  \resumeItem{Documented architecture, developer workflow, session management, provider setup, security mode, audit logs, and visible repository rules so the app remains a knowledge-base workbench rather than a black-box chatbot.}
\end{itemize}

\projectHeading{\href{https://github.com/Barytes/oh-share-it}{oh-share-it} -- Shared Research Context Layer}{MVP / in progress}
\begin{itemize}
  \resumeItem{Designed and implemented an early file-based context layer for research groups, with a \texttt{server + cli + client + tests} structure and explicit routing, URI, sync, rules, index, and JSON API surfaces.}
  \resumeItem{Separated the product boundary from \texttt{gogo}: \texttt{oh-share-it} owns shared knowledge governance, contribution/writeback workflows, routeable context, and agent-facing tools for team knowledge reuse.}
\end{itemize}

\projectHeading{context-core -- Eval-First Agent Context Infrastructure}{V0 design / skeleton}
\begin{itemize}
  \resumeItem{Specified a minimal context-infra core around \texttt{Source}, \texttt{ContextUnit}, \texttt{ContextBundle}, \texttt{RouteDecision}, \texttt{Trace}, \texttt{EvalCase}, and \texttt{WritebackCandidate}.}
  \resumeItem{Scoped V0 to make context selection, packaging, provenance, failure labels, and writeback candidates observable and regression-testable across agent workflows.}
\end{itemize}

\projectHeading{\href{https://github.com/Barytes/my-little-chating-agent}{my-little-chating-agent} / my-little-agent-loop -- Tool-Using Agent Harness}{Prototype / rebuild scope}
\begin{itemize}
  \resumeItem{Built an early FastAPI-based agent playground with tool use, web search, page reading, local Markdown/FAISS retrieval, and workflow scanning primitives.}
  \resumeItem{Narrowed the next rebuild to runtime ownership: session state, permission-gated tool execution, patch-based editing, JSONL traces, replayable tasks, and evaluator stubs aligned with explicit context bundles.}
\end{itemize}

\section{Research}

\resumeSubheading{Strategic User Offloading and Service Provider Pricing in Mobile Edge Computing}{ICASSP 2026}
{Accepted paper}{}
\begin{itemize}
  \resumeItem{Modeled strategic user behavior and provider pricing in mobile edge computing using game-theoretic and optimization methods.}
  \resumeItem{Transferable signal for agent systems: formal modeling, mechanism-design intuition, quantitative reasoning, and experience turning ambiguous system behavior into analyzable structure.}
\end{itemize}

\section{Writing and Knowledge Systems}
\begin{itemize}
  \resumeItem{Maintains a local markdown knowledge base with \texttt{raw/} evidence, maintained \texttt{wiki/} pages, reusable \texttt{frameworks/}, generated static browsing views, and repository-level agent instructions.}
  \resumeItem{Writes applied analyses on agent context infrastructure, harness design, public research knowledge bases, AI-native product boundaries, and career positioning for agent systems roles.}
\end{itemize}

\section{Education}

\resumeSubheading{Sun Yat-sen University}{Guangzhou, China}
{B.S. and M.S. in Software Engineering}{}

\section{Skills}
\begin{itemize}
  \resumeItem{\textbf{Agent systems:} LLM applications, RAG, tool use, agent loops, context routing, knowledge-base workflows, eval design, trace / replay planning, writeback workflows.}
  \resumeItem{\textbf{Engineering:} Python, TypeScript / JavaScript, FastAPI, Tauri desktop apps, CLI tools, REST / JSON APIs, Git, Markdown-based knowledge systems.}
  \resumeItem{\textbf{Research and product:} technical writing, system framing, product boundary definition, literature synthesis, game-theoretic modeling, optimization.}
\end{itemize}

\end{document}
